%%%%%%%%%%%%%%%%%%%%%%%%%%%%%%%%%%%%%%%%%%%%%%%
\begin{thebibliography}{900}

%\bibitem%[ja90]
%{ja90} J. Almeida, \emph{On pseudovarieties, varieties of languages, filters of congruences, pseudoidentities and related topics}. Algebra Universalis, \textbf{27} (1990), pp.~333--350.

\bibitem%[AntFLT26]
{AntFLT26} Anthropic, \emph{Formalizing Fermat's Last Theorem}. Anthropic Research, 4 September 2026. \url{https://www.anthropic.com/research/formalizing-fermats-last-theorem}.

\bibitem%[AntFLTRepo]
{AntFLTRepo} Anthropic, \emph{Fermat's Last Theorem in Lean}. Code repository. \url{https://github.com/anthropics/fermats-last-theorem}.

\bibitem%[AGMAI26]
{AGMAI26} Advisory Group on Mathematics and Artificial Intelligence, \emph{Responsible release of AI-generated mathematics}. 29 September 2026. \url{https://agmai.org/general-sep29/}.

\bibitem%[ABS00]
{ABS00}
F.~A. Al-Agl, R. Brown and R. Steiner, \emph{Multiple categories: the equivalence of a globular and a cubical approach}. arXiv preprint math/0007009 (2000).

\bibitem%[Lei26]
{Lei26} J. Alper, M. Barany, A. Chavarri Villarello, S. Dahmen, W. Dean, K. Ganapathy, M. Harris, D. Holmes, M. Jamnik, S. Kelk, B. Kra, U. Martin, B. Naskr\c{e}cki, R. Ochigame, J. Portegies and J. Schmitt, \emph{Leiden Declaration on Artificial Intelligence and Mathematics}. June 2026. arXiv preprint arXiv:2609.35334 (2026). \url{https://leidendeclaration.ai}.

\bibitem%[an83]
{an83} H. Andr\'{e}ka and I. N\'{e}meti, \emph{Generalization of the concept of variety and quasivariety to partial algebras through category theory.} Dissertationes Math. (Rozprawy Mat.) \textbf{204} (1983).

\bibitem%[abgmmm23]
{abgmmm23} D. Ara, A. Burroni, Y. Guiraud, P. Malbos, F. M\'{e}tayer and S. Mimram, \emph{Polygraphs: from rewriting to higher categories.} arXiv preprint arXiv:2312.00429 (2023).

\bibitem%[Ash26]
{Ash26} J. Asher, \emph{Formalize everything, now!} Blog post, 28 September 2026. \url{https://justinasher.me/formalize-everything-now}.

\bibitem%[BW85]
{BW85} M. Barr and Ch. Wells, \emph{Toposes, triples and theories}. Springer-Verlag, New York, 1985. 

\bibitem%[BW12]
{BW12} M. Barr and Ch. Wells, Category theory for computing science, Prentice-Hall International Series in Computer Science (1995); reprinted in: Reprints in Theory and Applications of Categories 22 (2012).


\bibitem%[jb67]
{jb67} J. B\'{e}nabou, \emph{Introduction to bicategories}. In Reports of the Midwest Category Seminar, volume 47 of Lecture
Notes in Mathematics, pp.~1--77. Springer, 1967.

\bibitem%[jb68]
{jb68} J. B\'{e}nabou, \emph{Structures alg\'{e}briques dans les cat\'{e}gories}. Cahiers Topologie G\'{e}om. Diff\'{e}rentielle \textbf{10} (1968),  pp.~1--126.

\bibitem{BN98}
F. Baader and T. Nipkow, \emph{Term rewriting and all that}. Cambridge University Press, 1999.

\bibitem{Ben70}
D.~B. Benson, \emph{Syntax and semantics: A categorical view}. Information and Control,
\textbf{17}, 1970, pp. 145-160.

\bibitem{Ben74}
D.~B. Benson, \emph{Semantic preserving translations}. Mathematical Systems Theory \textbf{8} (1974),  pp.~105--126.

\bibitem%[db75]
{db75} D.~B. Benson, \emph{The basic algebraic structures in categories of derivations}. Information and Control, \textbf{28} (1975), pp.~2--29.

\bibitem%[gb15]
{gb15} G.~M. Bergman, \emph{An invitation to general algebra and universal constructions}. Second edition. Universitext. Springer, Cham, 2015.

\bibitem%{BKV03}
{BKV03} M. Bezem, J.~W. Klop and R.~C. De Vrijer (Eds), \emph{Term rewriting systems. Terese}, Cambridge University Press, 2003.    

\bibitem%[bl05]
{bl05} T.~S. Blyth, \emph{Lattices and ordered algebraic structures}. Universitext. Springer-Verlag London, Ltd., London, 2005.

\bibitem%[arXiv26]
{arXiv26} K. Boboris, \emph{Fair moderation, equitable access, and AI: arXiv's updated rate limit policy}. arXiv blog, 1 October 2026. \url{https://blog.arxiv.org/2026/10/01/updated-rate-limit-policy/}.

\bibitem%[BO93]
{BO93} R.~V. Book and F. Otto, \emph{String-rewriting systems}.
Texts and Monographs in Computer Science. Springer-Verlag, New York, 1993. 

\bibitem%[nB70]
{nB70} N. Bourbaki, \emph{Th{\'{e}}orie des ensembles}. Hermann, Paris, 1970.

\bibitem%[DB17]
{DB17} D. Bourn, \emph{From groups to categorial algebra. Introduction to protomodular and Mal'tsev categories}. Compact Textbooks in Mathematics. Birkh\"{a}user/Springer, Cham, 2017.

%\bibitem%[Bra69]
%{Bra69} W.~S. Brainerd, \emph{Tree generating regular systems}. IC \textbf{14} (1969), pp.~217--231.

\bibitem%[BH81]
{BH81} R. Brown and P.~J. Higgins, \emph{The equivalence of $\infty $-groupoids and crossed complexes}. Cahiers de topologie et g\'{e}om\'{e}trie diff\'{e}rentielle, 22(4) (1981), pp.~371--386.

\bibitem%[rb89]
{rb89} J.~R. B\"{u}chi, \emph{Finite automata, their algebras and grammars. Towards a theory of formal expressions.} Edited and with a preface by Dirk Siefkes. Springer-Verlag, New York, 1989.

\bibitem%[bs67]
{bs67} P. Burmeister and J. Schmidt, \emph{On the completion of partial algebras}. Colloq. Math. \textbf{17} (1967), pp.~235--245.

\bibitem%[pb70]
{pb70} P. Burmeister, \emph{Free partial algebras}. J. Reine Angew. Math. \textbf{241} (1970), pp.~75--86.

\bibitem%[pb82]
{pb82} P. Burmeister, \emph{Partial algebras--survey of a unifying approach towards a twovalued
model theory for partial algebras}. Algebra Universalis, \textbf{15} (1982), pp.~306--358.

\bibitem%[pb86]
{pb86} P. Burmeister, \emph{A model theoretic oriented approach to partial algebras. Introduction to theory and application of partial algebras. Part I}. Mathematical Research, 32. Akademie-Verlag, Berlin, 1986.

\bibitem%[pbbw92a]
{pbw92a} P. Burmeister and B. Wojdy{\l}o, \emph{The meaning of basic category-theoretical notions in some categories of partial algebras. I}. Demonstratio Math. \textbf{25} (1992), pp.~583--602.

\bibitem%[pbw92b]
{pbw92b} P. Burmeister and B. Wojdy{\l}o, \emph{The meaning of basic category theoretical notions in some categories of partial algebras. II}. Products and coproducts. Demonstratio Math. \textbf{25} (1992), 
pp.~973--994.

\bibitem%[pb93]
{pb93}P. Burmeister, Partial algebras--{A}n introductory survey. Algebras and orders (Montreal, PQ, 1991), pp.~1--70, NATO Adv. Sci. Inst. Ser. C: Math. Phys. Sci., 389, Kluwer Acad. Publ., Dordrecht, 1993.

\bibitem%[pbw96]
{pbw96} P. Burmeister and B. Wojdy{\l}o, \emph{The meaning of basic category theoretical notions in some categories of partial algebras. III}. Demonstratio Math. \textbf{29} (1996), pp.~825--848.

\bibitem%[pb02]
{pb02} P. Burmeister, \emph{Lecture Notes on Universal Algebra--Many-Sorted Partial Algebras}. Fragment of lectures hold in the summer terms of 1994, 1996, 1998, 2000 and 2002 at the Darmstadt University of Technology. Download via http://www.mathematik.tu-darmstadt.de/{\raise.17ex\hbox{$\scriptstyle\mathtt{\sim}$}}burmeister/

\bibitem%[bs81]
{bs81} S. Burris and H.~P. Sankappanavar, \emph{A course in universal algebra}. Springer-Verlag, New York-Berlin, 1981.

\bibitem%[Buz26]
{Buz26} K. Buzzard, \emph{To grieve, or not to grieve?} Xena (blog), 1 October 2026. \url{https://xenaproject.wordpress.com/2026/10/01/to-grieve-or-not-to-grieve/}.

\bibitem{Bur91}
A. Burroni, \emph{Higher-dimensional word problem.} Category theory and computer science (Paris, 1991), pp.~ 94--105, Lecture Notes in Comput. Sci., \textbf{530}, Springer, Berlin, 1991.

\bibitem{Bur93}
A. Burroni, \emph{Higher-dimensional word problems with applications to equational logic}. Theoretical computer science \textbf{115} (1993), pp.~43-62.

\bibitem%[CH06]
{CH06} F. Cardone and J.~R. Hindley, \emph{History of lambda-calculus and combinatory logic}. In: D.M. Gabbay and J. Woods (Eds.) Handbook of the History of Logic.  Vol.~\textbf{5}. Logic from Russell to Church (2009), pp.~723--817. Elsevier/North-Holland, Amsterdam, 2009.

\bibitem%[CM82]
{CM82} B. Chandler and W. Magnus, \emph{The History of Combinatorial Group Theory}. Sources in the History of Mathematics and Physics Sciences \textbf{9}, 1982. 

\bibitem%[Chu32]
{ChuI32} A. Church, \emph{A set of postulates for the foundation of logic}. Ann. of Math. \textbf{33}(2) (1932), pp.~346--366.

\bibitem%[Chu32]
{ChuII32} A. Church, \emph{A set of postulates for the foundation of logic}. (Second paper.)  Ann. of Math. \textbf{34} (2) (1932), pp.~839--864.

\bibitem%[Chu36]
{Chu36} A. Church, \emph{An unsolvable problem of elementary number theory}. American Journal of Mathematics, \textbf{58} (2) (1936), pp.~345--363.

\bibitem%[Chu41]
{Chu41} A. Church, \emph{The calculi of lambda-conversion}. Ann. of Math. Studies \textbf{6}, Princeton N.~J., 1941.

\bibitem%[cgw95]
{cgw95} I. Cla{\ss}en, M. Gro{\ss}e-Rhode and U. Wolter, \emph{Categorical concepts for parameterized partial specifications}. ADT-COMPASS Workshop and TAPSOFT Conference (Caldes de Malavella, 1992; Orsay, 1993). Math. Structures Comput. Sci. \textbf{5} (1995), pp.~153--188.

\bibitem{CM15}
F. Clerc  and S. Mimram, \emph{Presenting a category modulo a rewriting system}. 26th International Conference on Rewriting Techniques and Applications (RTA' 2015). Editor: Maribel Fern\'{a}ndez; pp.~89--105.  Schloss Dagstuhl-Leibniz-Zentrum f\"ur Informatik, 2015.

\bibitem%[BCM20]
{BCM20} K. Buzzard, J. Commelin and P. Massot, \emph{Formalising perfectoid spaces}. In Proceedings of the 9th ACM SIGPLAN International Conference on Certified Programs and Proofs (CPP 2020), pp.~299--312. ACM, 2020.

% [No citada des del 04/10/2026, en escurçar la secció 7.5]
%\bibitem%[CIM26]
%{CIM26} M. Codreanu, A. Imas, J. Mateos-Garcia et al., \emph{AI in science: early insights}. Google, Google DeepMind and MIT FutureTech, September 2026. \url{https://ai.google/static/documents/AI-in-Science.pdf}.

% [No citada des del 04/10/2026, en escurçar la secció 7.5]
%\bibitem%[AIM26]
%{AIM26} M. Colbrook, \emph{AIM: an invitation to explore mathematics together}. Guest post on T.~Tao's blog, on behalf of S. Brunton, M. Colbrook, M.~V. de Hoop, G. Stepaniants, A. Townsend and R. Ward, 2 October 2026. \url{https://terrytao.wordpress.com/2026/10/02/aim-an-invitation-to-explore-mathematics-together/}.

\bibitem%[CLG25]
{CLG25} E. Cosme Ll\'{o}pez and L. Gong, \emph{An introduction to Lean~4}. Zenodo, 2025. \url{https://doi.org/10.5281/zenodo.17079867}.

\bibitem%[CosRepo]
{CosRepo} E. Cosme Ll\'{o}pez, \emph{Mslang: Lean~4 formalisation of Eilenberg theorems for many-sorted formations}. Code repository. \url{https://github.com/encosllo/MSLang}.

\bibitem%[CosArch]
{CosArch} E. Cosme Ll\'{o}pez, \emph{Architecture: concurrent mathematical development and formal verification (local-first)}. Specification, revision~6. \url{https://github.com/encosllo/MSLang/blob/main/Architecture.md}.

\bibitem%[CVST04]
{CVST04} J. Climent Vidal and J. Soliveres Tur,\emph{ On many-sorted algebraic closure operators}. Math. Nachr. \textbf{266} (2004), pp.~81--84. 

\bibitem%[CVST10a]
{CVST10a} J. Climent Vidal and J. Soliveres Tur,\emph{ Birkhoff-Frink representations as functors}. Math. Nachr. \textbf{283} (2010), pp.~686--703. 

\bibitem%[CVST10b]
{CVST10b} J. Climent Vidal and J. Soliveres Tur,\emph{ A 2-categorical framework for the syntax and semantics of many-sorted equational logic.} Rep. Math. Logic  \textbf{45} (2010), pp.~37--95.

\bibitem%[CVCL16]
{CVCL16} J. Climent Vidal and E. Cosme Ll\'{o}pez, \emph{Eilenberg theorems for many-sorted formations}.  Houston Journal of Mathematics, \textbf{45} (2019), pp.~321--369.

\bibitem%[CVCL18]
{CVCL18} J. Climent Vidal and E. Cosme Ll\'{o}pez, \emph{Congruence based proofs of the recognizability theorems for free many-sorted algebras}. Journal of Logic and Computation, \textbf{30} (2020), pp.~561--633.

\bibitem%[CVCL23]
{CVCL23} J. Climent Vidal and E. Cosme Ll\'{o}pez, \emph{Functoriality of the Schmidt Construction}. Logic Journal of the IGPL, \textbf{31}(5) (2023), pp.~822--893.

\bibitem%[CVCLRM25]
{CVCLRM25} J. Climent Vidal, E. Cosme Ll\'{o}pez and R. Ruiz Mora, \emph{Riguet and generalized congruences on a category. Relationships and applications}. Preprint (2025), arXiv:2505.15767.

\bibitem%[CVCLRM26]
{CVCLRM26}
Climent Vidal, J., Cosme Llópez, E., Ruiz Mora, R.: From higher-order rewriting systems to higher-order categorial algebras and higher-order Curry-Howard isomorphisms (2026), arXiv:2402.12051.

%\bibitem%[CS12]
%{CS12}
%J. Climent Vidal and J. Soliveres Tur, \emph{On the morphisms and transformations of Tsuyoshi Fujiwara (as a concretion of a bidimensional many-sorted general algebra and its application to the equivalence between many-sorted clones and algebraic theories)}, eprint arXiv:1201.5239, 2012.

\bibitem%[Coh81]
{pC81} P. Cohn, \emph{Universal algebra}, D. Reidel Publishing Company, 1981.

%\bibitem%[Coq86]
%{Coq86} T. Coquand and G. Huet, \emph{The calculus of constructions}. Information and Control, \textbf{76} (1988), pp.~95--120.

\bibitem%[CurI30]
{CurI30} H.~B. Curry, \emph{Grundlagen der Kombinatorischen Logik. Teil I}. American Journal of Mathematics, \textbf{52} (3) (1930), pp.~509--536.

\bibitem%[CurII30]
{CurII30} H.~B. Curry, \emph{Grundlagen der Kombinatorischen Logik. Teil II}. American Journal of Mathematics, \textbf{52} (4) (1930), pp.~789--834.

\bibitem%[CF58]
{CF58} H.~B. Curry and R. Feys, \emph{Combinatory Logic, vol. $\mathrm{I}$}. Studies in Logic and the Foundations of Mathematics. North-Holland Publishing Co., Amsterdam, 1958.

%\bibitem%[Tata08]
%{Tata08} H. Comon, M. Dauchet, R. Gilleron, F. Jacquemard, D. Lugiez, S. Tison, C. L\"{o}ding, and M. Tommasi,
%\emph{Tree automata techniques and Applications}. Available on: {\tt http://tata.gforge.inria.fr}, 2008.

%\bibitem%[Cou89]
%{Cou89} B. Courcelle, \emph{On recognizable sets and tree automata}. In H. A\"{\i}t-Kaci and M. Nivat, editors. \emph{Algebraic Techniques: Resolution of Equations in Algebraic Structures}, Chapter 3, pp.~93--126, Academic press, Inc., 1989.

%\bibitem%[Cou96]
%{Cou96} B. Courcelle, \emph{Basic notions of universal algebra for language theory and graph grammars}, Theoret. Comput. Sci. \textbf{163} (1996), pp.~1--54.
%

\bibitem%[Die26]
{Die26} T.~G. Dietterich, Post on X, 13 September 2026. \url{https://x.com/tdietterich/status/2099197831580090622}.

\bibitem%[DP02]
{DP02} B.~A. Davey and H.~A. Priestley, \emph{Introduction to lattices and order}. Second edition. Cambridge University Press, New York, 2002.

\bibitem%[dMU21]
{dMU21} L. de Moura and S. Ullrich, \emph{The Lean~4 theorem prover and programming language}. In Automated Deduction -- CADE~28, Lecture Notes in Computer Science \textbf{12699}, pp.~625--635. Springer, 2021.

\bibitem%[D82]
{D82} M. Davis, \emph{Computability and unsolvability}. (Enl. version of the orig. publ. McGraw- Hill Book Company, New York 1958). Dover Publications, Inc. 
New York, 1982.

\bibitem%[D04]
{D04} M. Davis (ed.), \emph{The undecidable. Basic papers on undecidable propositions, unsolvable problems and computable functions}. Corrected reprint of the 1965 original. Dover Publications, Inc., Mineola, New York, 2004.

%\bibitem%[d63]
%{d63} R. Dedekind, \emph{Essays on the theory of numbers. I: Continuity and irrational numbers. II: The nature and meaning of numbers}. Authorized translation by Wooster Woodruff Beman Dover Publications, Inc., New York, 1963.

\bibitem%[MD11]
{MD11} M. Dehn, \emph{\"{U}ber unendliche diskontinuierliche Gruppen}, Mathematische Annalen \textbf{71} (1911), pp.~116--144.

%\bibitem%[Don65]
%{Don65} J.~E. Doner, \emph{Decidability of the weak second-order theory of two successors}. Notices Amer. Math. Soc. \textbf{12} (1965), Abstract 65T-468, p.~819.
%
%\bibitem%[Don70]
%{Don70} J.~E. Doner, \emph{Tree acceptors and some of their applications}. J. CSS \textbf{4} (1970), pp.~406--451.

\bibitem{Ehr64}
C. Ehresmann, \emph{Categories et structures: extraits}. Seminaire Ehresmann. Topologie et geometrie differentielle \textbf{6} (1964) pp.~1-31.

\bibitem%[ekkk74]
{ekkk74} H. Ehrig, K.-D. Kiermeier, H.-J. Kreowski, and W. K\"{u}hnel,\emph{ Universal theory of automata. A categorical approach.} Teubner Studienb\"{u}cher: Informatik. B.~G. Teubner, Stuttgart, 1974.

\bibitem%[EM45]
{EM45} S. Eilenberg and S. MacLane, \emph{General theory of natural equivalences}. Trans. Amer. Math. Soc. \textbf{58} (1945), pp.~231--294. 

\bibitem%[EW67]
{EW67} S. Eilenberg and J.~B. Wright, \emph{Automata in general algebras}. Information and Control \textbf{11} (1967),  pp.~452--470.

%\bibitem%[Eng75]
%{Eng75} J. Engelfriet, \emph{Bottom-up and top-down tree transformations. A comparison.} MST \textbf{9} (1957), pp.~198--231.

%\bibitem%[EngSch77]
%{EngSch77} J. Engelfriet and E.~M. Schmidt, \emph{IO and OI}. I. J. Comput. System Sci. \textbf{15}  (1977), pp.~328--353.

%\bibitem%[je15]
%{je15} J. Engelfriet, \emph{Tree automata and tree grammars}. Arxiv preprint arXiv:1510.02036v1, 2015.

%\bibitem%[Fuj59]
%{Fuj59}
%T. Fujiwara, \emph{On mappings between algebraic systems}, Osaka Math. J., \textbf{11}, pp.~153--172, 1959.

%\bibitem%[Fuj59]
%{Fuj60}
%T. Fujiwara, \emph{On mappings between algebraic systems II}, Osaka Math. J., \textbf{12}, pp.~253--268, 1960.

\bibitem%[end77]
{end77} H.~B. Enderton, \emph{Elements of set theory}. Academic Press [Harcourt Brace Jovanovich, Publishers], New York-London, 1977. 

\bibitem%[Formalpedia26]
{Formalpedia26} Prove2me, \emph{Formalpedia}. \url{https://prove2.me/formalpedia}. Accessed September 2026.

\bibitem%[Gar26]
{Gar26} D.~Z. Garza, Post on X, 10 September 2026. \url{https://x.com/dzackgarza/status/2098092343702548680}.

\bibitem%[GJ17]
{GJ17} N. Gambino and A. Joyal, \emph{On operads, bimodules and analytic functors}. Mem. Amer. Math. Soc.  \textbf{249} (2017), no. 1184.

\bibitem%[GK15]
{GK15} C. Ganesamoorthy and S.~G. Karpagavalli, \emph{A congruence relation in partially ordered sets}. Thai J. Math. \textbf{13} (2015), pp.~399-403.

\bibitem%[GRMSVCL26]
{GRMSVCL26} L. Gong, R. Ruiz Mora, N. Sanmart\'{\i}n Vich and E. Cosme Ll\'{o}pez, \emph{A Kleene theorem for free many-sorted algebras}. arXiv preprint arXiv:2606.29939 (2026).

\bibitem%[GS84]
{GS84} F. G\'{e}cseg and M. Steinby, \emph{Tree automata}. Akad\'{e}miai Kiad\'{o}, Budapest, 1984.

% [No citada des del 04/10/2026, en escurçar la secció 7.5]
%\bibitem%[Gut26]
%{Gut26} L. Guth, \emph{What math means to me}. arXiv preprint arXiv:2609.32028 (2026).

%\bibitem%[GS97]
%{GS97} F. G\'{e}cseg and M. Steinby, \emph{Tree languages}. In G. Rozenberg and A. Salomaa, editors.
%\emph{Handbook of formal languages. Vol. 3. Beyond words}, Chapter 1, pp.~1--68, Springer-Verlag, Berlin, 1997.

\bibitem%[G31]
{G31} K. G\"{o}del, \emph{\"{U}ber formal unentscheidbare S\"{a}tze der Principia Mathematica und verwandter Systeme. I.} 
Monatsh. f. Mathematik und Physik \textbf{38} (1931), pp.~173--198.

\bibitem%[GTW]
{gtw85} J. Goguen, J. Thatcher, and E. Wagner, \emph{An initial algebra approach to the specification, correctness, and implementation of abastract data types}. In R. Yeh, editor. \emph{Current Trends in Programming Methodology, IV}, Chapter 5, pp.~80--149, Prentice Hall, New Jersey, 1978.

\bibitem%[GM85]
{gm85} J. Goguen and J. Meseguer, \emph{Completeness of many-sorted equational logic}. Houston Journal of Mathematics, \textbf{11} (1985), pp.~307--334.

%\bibitem%[GB86]
%{gb86} J. Goguen and R. Burstall, \emph{A study in the foundations of programming methodology: Specifications, institutions, charters and parchments}. In D. Pitt et alii. eds. \emph{Proc. Summer Workshop on Category Theory and Computer Programming},
%Springer-Verlag (1986), pp.~313--333.

%\bibitem%[GW23]
%{GW23} U. G\"{o}rtz and T. Wedhorn, \emph{Algebraic geometry II: Cohomology of Schemes}. Springer Stud. Math--Master. Springer Spektrum, Wiesbaden, 2023.

\bibitem%[Grt08]
{gG08} G. Gr\"{a}tzer, \emph{Universal algebra. With appendices by Gr\"{a}tzer, Bjarni J\'{o}nsson, Walter Taylor, Robert W. Quackenbush, G\"{u}nter H. Wenzel, and Gr\"{a}tzer and W.~A. Lampe}. Revised reprint of the 1979 second edition. Springer, New York, 2008.

%\bibitem%[Gry66]
%{jwG66} J.~W. Gray, \emph{Fibred and cofibred categories.} In S. Eilenberg et alii. eds. \emph{Proc. Conf. Categorical Algebra}, Springer-Verlag (1966), pp.~21--83.

%\bibitem%[gree52]
%{gree52} J.~A. Green, \emph{A duality in abstract algebra}. J. London Math. Soc., \textbf{27} (1952), pp.~64--73.

\bibitem%[Gro71]
{Gro71} A. Grothendieck. \emph{Cat\'{e}gories fibr\'{e}es et
descente (Expos\'{e} VI)}.  In A. Grothendieck, editor,
Rev\^{e}tements \'{e}tales et groupe fondamental (SGA 1),
Springer-Verlag, Berlin $\cdot$ Heidelberg $\cdot$ New York, 1971, pp.~145--194.

\bibitem{Gui04}
Y. Guiraud, \emph{Pr\'esentations d'op\'erades et syst\`emes de r\'e\'ecriture}. PhD thesis, Universit\'e Montpellier 2, 2004.

\bibitem{GM09}
Y. Guiraud and P. Malbos, \emph{Higher-dimensional categories with finite derivation type}. Theory and Applications of Categories Vol. \textbf{22} (2009), pp.~420--478.

\bibitem{Gui19}
Y. Guiraud, \emph{Rewriting methods in higher algebra}, Mémoire d'Habilitation à Diriger des Recherches (2019), Univ. Paris 7.

\bibitem%[Har26]
{Har26} K. Hartnett, \emph{The proof in the code: how a truth machine is transforming math and AI}. Farrar, Straus and Giroux, New York, 2026.

\bibitem%[H05]
{H05} E. Harzheim, \emph{Ordered sets}. Advances in Mathematics (Springer), \textbf{7}. Springer, New York, 2005. 

\bibitem%[Her30]
{Her30} J. Herbrand, \emph{Recherches sur la th\'{e}orie de la d\'{e}monstration}, PhD Thesis, Universit\'{e} de Paris, 1930.

\bibitem%[hs73]
{hs73} H. Herrlich and G.~E. Strecker, \emph{Category theory: an introduction}. Allyn and Bacon Series in Advanced Mathematics. Allyn and Bacon Inc., Boston, Mass., 1973.


\bibitem%[Hig63]
{Hig63} P.~J. Higgins, \emph{Algebras with a scheme of operators}. Math. Nachr. \textbf{27} (1963), pp.~115--132. 

\bibitem%[h71]
{h71} P.~H. Higgins, \emph{Notes on categories and groupoids}. Van Nostrand Reinhold Co., London-New York-Melbourne, 1971.

\bibitem%[hs07]
{hs07} H.-J. Hoehnke and J. Schreckenberger, \emph{Partial algebras and their theories}. Shaker Verlag, Aachen, 2007. 

%\bibitem%[HoTT13]
%{HoTT13} \emph{Homotopy Type Theory: Univalent Foundations of Mathematics. The Univalent Foundations Program}, Institute for Advanced Study, 2013.

\bibitem%[Hook85]
{Hook85} J.~L. Hook, \emph{A note on interpretations of many-sorted theories}. The Journal of Symbolic Logic, Vol. \textbf{50} (1985), pp.~372--374.

\bibitem%[Hotz65]
{Hotz65} G. Hotz, \emph{Eine Algebraisierung des Syntheseproblems von Schaltkreisen. $\mathrm{I}$.} 
Elektronisch Informationsverarbeitung und Kybernetik, Vol. \textbf{1} (1965), pp.~185--205.

\bibitem%[Hotz65a]
{Hotz65a} G. Hotz, \emph{Eine Algebraisierung des Syntheseproblems von Schaltkreisen. $\mathrm{II}$.} 
Elektronisch Informationsverarbeitung und Kybernetik, Vol. \textbf{1} (1965), pp.~209--231.

\bibitem%[Hotz66]
{Hotz66} G. Hotz, \emph{Eindeutigkeit and Mehrdeutigkeit formaler Sprachen.}
Elektronisch Informationsverarbeitung und Kybernetik, Vol. \textbf{2} (1966), pp.~235--246

\bibitem%[H80]
{H80} W.~A. Howard, \emph{The formulae-as-types notion of construction.} In J.~P. Seldin and J.~R. Hindley, editors. \emph{To H.~B. Curry: Essays on Combinatory Logic, Lambda Calculus and Formalism}, pp.~479--490, Academic Press, London-New York, 1980.

\bibitem%[bj99]
{bj99} B. Jacobs, \emph{Categorical logic and type theory}. Stud. Logic Found. Math., \textbf{141}. North-Holland Publishing Co., Amsterdam, 1999.

\bibitem%[MJ88]
{MJ88} M. Jantzen,  
\emph{Confluent string rewriting}.
EATCS Monographs on Theoretical Computer Science, \textbf{14}. Springer-Verlag, Berlin, 1988.

\bibitem%[J87]
{J87} M. Johnson, \emph{Pasting diagrams in n-categories with applications to coherence theorems and categories of paths}, PhD thesis, University of Sydney, 1987.

\bibitem%[JY21]
{JY21} N. Johnson and D. Yau, \emph{2-dimensional categories}. Oxford University Press, Oxford, 2021.

%\bibitem%[jt51]
%{jt51} B. J\'{o}nsson and A. Tarski, \emph{Boolean algebras with operators, I}. Amer. J. Math., \textbf{73} (1951), pp.~891--939.
%
%\bibitem%[jt52]
%{jt52} B. J\'{o}nsson and A. Tarski, \emph{Boolean algebras with operators, II}. Amer. J. Math., \textbf{74} (1952), pp.~127--162.

\bibitem%[Jon72]
{bJ72}
B. J\'{o}nsson, \emph{Topics in universal algebra}, Springer-Verlag, 1972.

%\bibitem%[kl80]
%{kl80} M. Kline, \emph{Mathematics. The loss of certainty}. Oxford University Press, New York, 1980.

% [No citada des del 04/10/2026, en escurçar la secció 7.5]
%\bibitem%[Kal26]
%{Kal26} G. Kalai, \emph{``Be a Grothendieck!'': on AI and mathematics}. Combinatorics and more (blog), 22 September 2026. \url{https://gilkalai.wordpress.com/2026/09/22/be-a-grothendieck-on-ai-and-mathematics/}.

\bibitem%[kr72]
{kr72} H. Kaphengst and H. Reichel, \emph{Operative Theorien und Kategorien von operativen Systemen}. In  \emph{Studien zur Algebra und ihre Anwendungen (mit Anwendungen in der Mathematik, Physik und Rechentechnik)}, pp.~41--56. Schr. Zentralinst. Math. Mech. Akad. Wissensch. DDR, No. 16. Akademie Verlag, Berlin, 1972.

\bibitem%[kkm00]
{kkm00} M. Kilp, U. Knauer and A.~V. Mikhalev, \emph{Monoids, acts and categories. With applications to wreath products and graphs. A handbook for students and researchers}. Walter de Gruyter \!\!$\And$\!\! Co., Berlin, 2000.

\bibitem{Klo92}
J.~W. Klop, \emph{Term rewriting systems}. Handbook of Logic in Computer Science, Vol. \textbf{2}, Oxford University Press, 1992, pp.~1-117.

\bibitem%[agk67]
{agk67} A.~G. Kurosh, \emph{Alg\`{e}bre g\'{e}n\'{e}rale}. (French) Traduit du Russe par J.-P. Peaudecerf avec la participation de G. Lallement Collection Universitaire de Math\'{e}matiques, No. 22. Dunod, Paris, 1967.

\bibitem%[Kou26]
{Kou26} D. Koukoulopoulos, \emph{A CERN for AI-assisted science?} Guest post on T.~Tao's blog, 17 September 2026. \url{https://terrytao.wordpress.com/2026/09/17/a-cern-for-ai-assisted-science/}.

\bibitem{Laf03}
Y. Lafont, \emph{Towards an algebraic theory of Boolean circuits}. Journal of Pure and Applied Algebra \textbf{184} (2003), pp.~257-310.

\bibitem{LMW10}
Y. Lafont, F. M\'{e}tayer and K. Worytkiewicz, \emph{A folk model structure on omega-cat}. Advances in Mathematics \textbf{224}.3 (2010), pp.~1183--1231.

\bibitem%[Lk81]
{Lk81} J. Lambek, \emph{The Influence of Heraclitus on Modern Mathematics}. In: Agassi, J., Cohen, R.~S. (eds) Scientific Philosophy Today. Boston Studies in the Philosophy of Science, vol \textbf{67}. Springer, Dordrecht, 1981, pp.~111--121.

\bibitem%[LS86]
{LS86} J. Lambek and P.~J. Scott, \emph{Introduction to higher order categorical logic}. Cambridge Stud. Adv. Math., \textbf{7}. Cambridge University Press, Cambridge, 1986.

\bibitem{Lan21}
M. Land, \emph{Introduction to Infinity-categories}. Springer Nature, 2021.

%\bibitem%[l98]
%{l98} M. V. Lawson, \emph{Inverse semigroups. The theory of partial symmetries}. World Scientific Publishing Co., Inc., River Edge, NJ, 1998. 

\bibitem{Law03}
F.~W. Lawvere, \emph{Functorial semantics of algebraic theories}. Proceedings of the National Academy of Sciences \textbf{50} (1963), pp.~869-872.

%\bibitem%[L94]
%{L94} F. W. Lawvere, \emph{Cohesive toposes and Cantor's ``lauter Einsen''}. Categories in the foundations of mathematics and language. Philos. Math. (3) \textbf{2} (1994), pp.~5--15. 

\bibitem{Lei01}
T. Leinster, \emph{A survey of definitions of n-category}. arXiv preprint math/0107188 (2001).

\bibitem%[L02]
{L02} T. Leinster, \emph{A survey of definitions of $n$-category}. Theory Appl. Categ. \textbf{10} (2002), pp.~1--70.

\bibitem{Lei04}
T. Leinster, \emph{Higher operads, higher categories}. No. \textbf{298}. Cambridge University Press, 2004.

\bibitem%[LG17]
{LG17} S. Lucas and R. Guti\'{e}rrez, \emph{Automatic Synthesis of Logical Models for Order-Sorted First-Order Theories}. Journal of Automated Reasoning, \textbf{59} (2017), pp.~381--412. DOI: 10.1007/s10817-017-9419-3.

%\bibitem%[Lall79]
%{Lall79} G. Lallement, \emph{Semigroups and combinatorial applications}. Pure and Applied Mathematics. A Wiley-Interscience Publication. John Wiley $\!\!\And\!\!$ Sons, New York-Chichester-Brisbane, 1979.

%\bibitem%[sl94]
%{sl94} S. Lang, \emph{Algebraic number theory}. Second edition. Graduate Texts in Mathematics, 110. Springer-Verlag, New York, 1994.

%\bibitem%[sl96]
%{sl96} S. Lang, \emph{Topics in cohomology of groups.} Translated from the 1967 French original by the author. Chapter X based on letters written by John Tate. Springer-Verlag, Berlin, 1996.

%\bibitem%[law96]
%{law96} F. W. Lawvere,
%\emph{Adjoints in and among bicategories}. In \emph{Logic and algebra}.
%Papers from the International Conference in memory of Roberto Magari held in Pontignano, April 26-30, 1994. Edited by Aldo Ursini and Paolo Aglian\`{o}. Lecture Notes in Pure and Applied Mathematics, 180, pp.~181--189. Marcel Dekker, Inc., New York, 1996.
%
%\bibitem%[lc63]
%{lc63} L. Lesieur and R. Croisot, \emph{Alg\`{e}bre noeth\'{e}rienne non commutative}. 
%Gauthier-Villars \!$\And$\! Cie, \'{E}diteur-Imprimeur-Libraire, Paris, 1963. 

\bibitem%[Lax26]
{Lax26} \'{E}. Bonnet, J. Dreier and C. Kuske, \emph{Lax: the Lean archive}. \url{https://laxarchive.org}. Accessed September 2026.

\bibitem%[LaxAbout26]
{LaxAbout26} \'{E}. Bonnet, J. Dreier and C. Kuske, \emph{About Lax}. \url{https://laxarchive.org/about.html}. Accessed September 2026.

\bibitem%[LeanFRO26]
{LeanFRO26} Lean FRO, \emph{About the Lean FRO}. Accessed 4 October 2026. \url{https://lean-lang.org/fro/about/}.

\bibitem%[Lit26]
{Lit26} D. Litt, \emph{A beginning for mathematics}. Proofs and Prompts (blog), 14 September 2026. \url{https://proofsandprompts.com/2026/09/14/a-beginning-for-mathematics/}.

\bibitem%[KR35]
{KR35} S.~C. Kleene and J.~B. Rosser, \emph{The inconsistency of certain formal logics}. Ann. of Math. (2) \textbf{36}(3) (1935), pp.~630--636.

\bibitem%[lj97]
{lj97} E.~S. Ljapin and A.~E. Evseev,
\emph{The theory of partial algebraic operations}. Translated from the 1991 Russian original and with a preface by J. M. Cole. Revised by the authors. Mathematics and its Applications, 414. Kluwer Academic Publishers Group, Dordrecht, 1997. 

\bibitem%[mathlib20]
{mathlib20} The mathlib Community, \emph{The Lean mathematical library}. In Proceedings of the 9th ACM SIGPLAN International Conference on Certified Programs and Proofs (CPP 2020), pp.~367--381. ACM, 2020.

\bibitem%[mathlibDocs26]
{mathlibDocs26} The mathlib Community, \emph{Mathlib documentation}. \url{https://leanprover-community.github.io/mathlib4_docs/}. Accessed September 2026.

% [No citada des del 04/10/2026, en escurçar la secció 7.5]
%\bibitem%[MAIT26]
%{MAIT26} \emph{Mathematics and AI: an open letter on human--AI collaboration in mathematics}. Open letter, accessed 4 October 2026. \url{https://sites.google.com/view/math-and-ai-together/home}.

\bibitem%[Mac98]
{sM98} S. Mac Lane, \emph{Categories for the working mathematician.} 2nd ed. Springer-Verlag, New York, 1998.

%\bibitem{Mac13}
%S. Mac Lane, \emph{Categories for the working mathematician}. Vol. 5. Springer Science \& Business Media, 2013.

%\bibitem%[MM69]
%{MM69} M. Magidor and G. Moran, \emph{Finite automata over finite trees}. Technical report, \textbf{30}, Hebrew University, Jerusalem, 1969.

%\bibitem%[Maib74]
%{Maib74} T. S. E. Maibaum, \emph{A generalized approach to formal languages}. J. Comput. System Sci.  \textbf{8} (1974), pp.~409--439.
%
%\bibitem%[Maib77]
%{Maib77} T. S. E. Maibaum, \emph{Erratum}. J. Comput. System Sci.  \textbf{14} (1977), p.~369.

\bibitem%[em76]
{em76} E.~G. Manes, \emph{Algebraic theories}, Springer-Verlag, New York, 1976.

\bibitem%[Mar47]
{Mar47} A. Markov, \emph{On the impossibility of certain algorithms in the theory of associative systems}. C. R. (Doklady) Acad. Sci. URSS (N.S.) \textbf{55} (1947), pp.~583--586.

%\bibitem%[LS84]
%{LS84} P. Martin-L{\"o}f, \emph{Intuitionistic type theory}, \textbf{9} Bibliopolis, Naples, 1984.

\bibitem%[m76]
{m76} G. Matthiessen, \emph{Theorie der {H}eterogenen {A}lgebren}, Mathematik-{A}rbeitspapiere, Nr. 3, Universit\"{a}t Bremen Teil A, Mathematische Forchungspapiere, 1976.

%\bibitem%[MY60]
%{MY60} R. McNaughton and H. Yamada, \emph{Regular expressions and state graphs for automata}. IRE Transactions on Electronic Computers, \textbf{EC-9} (1960), pp.~39--47.

\bibitem{Met03}
F. M\'{e}tayer, \emph{Resolutions by polygraphs}. Theory and Applications of Categories \textbf{11} (2003), pp.~148--184.

%\bibitem%{MW67}
%{MW67} J. Mezei and J. Wright, \emph{Algebraic automata and context-free sets}. Information and Control, \textbf{11} (1967),   pp.~3--29.

% [No citada des del 04/10/2026, en escurçar la secció 7.5]
%\bibitem%[MDZ26]
%{MDZ26} J. Miao, J.~R. Davis, Y. Zhang, J.~K. Pritchard and J. Zou, \emph{Reimagining research papers as interactive and reliable AI agents}. Nature (2026). \url{https://www.nature.com/articles/s41586-026-11044-y}.

\bibitem%[Mi14]
{Mi14} C.~F. Miller\! III, \emph{Turing machines to word problems.} In R. Downey, editor. \emph{Turing's legacy: developments from Turing's ideas in logic,} pp.~329--385, Lect. Notes Log., \textbf{42} Association for Symbolic Logic, La Jolla, CA; Cambridge University Press, Cambridge, 2014.

\bibitem{Mim14}
S. Mimram, \emph{Towards 3-dimensional rewriting theory}. Logical Methods in Computer Science \textbf{10} (2014), pp.~1--47.

%\bibitem%[Myh57]
%{Myh57} J. Myhill, \emph{Finite automata and the representation of events.}  WADD TR-57-624, Wright Patterson AFB, Ohio (1957), pp.~112--137.
%

\bibitem%[Mul24]
{Mul24} S. M\"{u}ller-Stach, \emph{Max Dehn, Axel Thue, and the Undecidable} In J. Lorenat, J. McClearly, V.~R. Remmert, D.~E. Rowe \& M. Senechal (Eds) \emph{Max Dehn: Polyphonic portrait}, Chapter 8, pp.~181--196, American Mathematical Society, \textbf{46}, 2024.

%\bibitem%[Ner58]
%{Ner58} A. Nerode, \emph{Linear automata transformation.} In Proc. American Mathematical Society, \textbf{9} (1958), pp.~541--544.

\bibitem%[N55]
{N55}
P.~S. Novikov, \emph{On the algorithmic unsolvability of the word problem in group theory}. Akademiya Nauk S.S.S.R. Matematicheskii Institut Trudy, \textbf{44}, Moscow (1955), pp.~1--143.

%\bibitem%[po11]
%{po11} S.~M. Pinto and M.~T. Oliveira-Martins, \emph{Congruences and ideals on Boolean modules: a heterogeneous point of view}. Math. Log. Quart. \textbf{57} (2011), pp.~571--581.
%
%\bibitem%[po12]
%{po12} S.~M. Pinto and M.~T. Oliveira-Martins, \emph{Congruences and ideals on Peirce algebras: a heterogeneous/homogeneous point of view}. Math. Log. Quart. \textbf{58} (2012), pp.~252--262.

\bibitem%[bp72]
{bp72} B.~I. Plotkin, \emph{Groups of automorphisms of algebraic systems}. Translated from the Russian by K.~A. Hirsch. Wolters-Noordhoff Publishing, Groningen, 1972.

\bibitem%[OAINS26]
{OAINS26} OpenAI, \emph{On the Navier--Stokes Millennium Prize Problem}. OpenAI, 8 September 2026. \url{https://openai.com/index/navier-stokes-solution/}.

\bibitem%[OAINSRepo]
{OAINSRepo} OpenAI, \emph{NavierStokesAndEuler: Lean certificates accompanying Navier--Stokes and Euler results}. Code repository. \url{https://github.com/openai/NavierStokesAndEuler}.

\bibitem%[OAIAG26]
{OAIAG26} OpenAI, \emph{Advisory Group on Mathematics and Artificial Intelligence}. OpenAI, 21 September 2026. \url{https://openai.com/index/advisory-group-on-mathematics-and-ai/}.

\bibitem%[PT91]
{PT91}H.-E. Porst and W. Tholen, \emph{Concrete dualities}. Category theory at work (Bremen, 1990), pp.~111--136. Res. Exp. Math., \textbf{18} Heldermann Verlag, Berlin, 1991.


\bibitem%[Poz26]
{Poz26} B. Pozsgay, Post on X, 23 September 2026. \url{https://x.com/pozsgaybalazs/status/2102896555418177667}.

\bibitem%[Post43]
{Post43} E.~L. Post, \emph{Formal reductions of the general combinatorial decision problem}. Amer. J. Math.  \textbf{65} (1943), pp.~197--215.

\bibitem%[Post44]
{Post44} E.~L. Post. Recursively enumerable sets of positive integers and their decision problems. Bulletin of the American Mathematical Society \textbf{50} (1944), pp.~284--316.

\bibitem%[Post47]
{Post47} E.~L. Post, \emph{Recursive unsolvability of a problem of Thue}. The Journal of Symbolic Logic \textbf{12}(1) (1947), pp.~1--11.

\bibitem%[Post65]
{Post65}  E.~L. Post, \emph{Absolutely unsolvable problems and relatively undecidable
propositions---account of an Anticipation}. In \emph{The undecidable. Basic papers on undecidable propositions, unsolvable problems and computable functions}, edited by Martin Davis, Raven Press Books, Ltd., Hewlett, N.~Y. 1965, pp.~340--433.

\bibitem%[Post04]
{Post04}  E.~L. Post, \emph{Absolutely unsolvable problems and relatively undecidable propositions---account of an Anticipation}. In \emph{The undecidable. Basic papers on undecidable propositions, unsolvable problems and computable functions}, edited by Martin Davis, Dover Publications, Inc., Mineola, New York, 2004, pp.~340--406.

\bibitem{Pow91}
A.~J. Power, \emph{An $n$-categorical pasting theorem}. In \emph{Category theory}, volume 1488 of Lecture Notes in Mathematics, pp.~326-358, Springer, 1991.

\bibitem{Pow89}
A.~J. Power, \emph{An abstract formulation for rewrite systems}. In \emph{Category Theory and Computer Science}, Manchester, UK, September 5--8, 1989 Proceedings, pp.~300--312, Springer Berlin Heidelberg, 2005.

\bibitem{Pow14}
J.~F. Power, \emph{An Early Completion Algorithm: Thue's 1914 Paper on the Transformation of Symbol Sequences}. In: A. Beckmann, E. Csuhaj-Varj\'{u} and K. Meer  (Eds.) Language, Life, Limits. CiE 2014. Lecture Notes in Computer Science, 
vol.~\textbf{8493},  pp.~343--346, Springer, Cham, 2014.

%\bibitem%[Rab69]
%{Rab69} M.~O. Rabin, \emph{Decidability of second-order theories and automata on infinite trees}. Trans. Amer. Math. Soc., \textbf{141} (1969), pp.~1--35.

%\bibitem%[RS59]
%{RS59} M.~O. Rabin and D. Scott, \emph{Finite automata and their decision problems}. IBM J. Res. Develop. \textbf{3} (1959) pp.~114--125.
%
%\bibitem%[rs04]
%{rs04} M. Raussen and Chr. Skau, \emph{Interview with Michael Atiyah and Isadore Singer}. Eur. Math. Soc. Newsl. \textbf{53}  (2004), pp.~24--30.

\bibitem%[P2M26]
{P2M26} Prove2me, \emph{Prove2me: humans and AI agents proving theorems in Lean~4}. \url{https://prove2.me}. Accessed September 2026.

\bibitem%[P2MStart26]
{P2MStart26} Prove2me, \emph{Getting started}. \url{https://prove2.me/start.md}. Accessed September 2026.

\bibitem%[P2MEil26]
{P2MEil26} Prove2me, \emph{Mission: Eilenberg theorems for many-sorted formations}. Accepted 9 September 2026. \url{https://prove2.me/missions/Eilenberg_Theorems_for_Many-Sorted_Formations}.

\bibitem%[P2MKle26]
{P2MKle26} Prove2me, \emph{Mission: A Kleene theorem for free many-sorted algebras}. Accepted 9 September 2026. \url{https://prove2.me/missions/A_Kleene_Theorem_for_Free_Many-Sorted_Algebras}.

\bibitem%[QW26]
{QW26} Y. Qi and M. Weber, \emph{Choir: an open protocol for distributed multi-agent autoformalization}. arXiv preprint arXiv:2609.31903 (2026).

\bibitem%[Rei87]
{Rei87} H.~Reichel, \emph{Initial Computability, Algebraic Specifications, and Partial Algebras}. International Series of Monographs on Computer Science, Clarendon Press, Oxford, 1987.

\bibitem%[r87]
{r87} H. Reichel, \emph{Initial computability, algebraic specifications, and partial algebras}. The Clarendon Press, Oxford University Press, New York, 1987.

\bibitem{RV16}
E. Riehl and D. Verity, \emph{Infinity category theory from scratch}. arXiv preprint arXiv:1608.05314 (2016).

\bibitem%[rv22]
{rv22} E. Riehl and D. Verity, \emph{Elements of $\infty$-category theory}. Cambridge University Press, Cambridge, 2022.

\bibitem%[R08]
{R08} S. Roman, \emph{Lattices and ordered sets}. Springer, New York, 2008.

\bibitem%[pr18]
{pr18} P. Roquette, \emph{The Riemann Hypothesis in Characteristic $p$ in Historical Perspective}, Springer Nature Switzerland AG, 2018.

%\bibitem%[kr96]
%{kr96} K.~I. Rosenthal, \emph{The theory of quantaloids}. Longman, Harlow, 1996.

%\bibitem%[cr10]
%{cr10} C. Russo,\emph{ Quantale modules and their operators, with applications}. J. Logic Comput. \textbf{20} (2010), pp.~917--946. 

\bibitem%[Sah26]
{Sah26} A. Sahai, \emph{We're gonna need a lot more mathematicians}. Guest post on T.~Tao's blog, 24 September 2026. \url{https://terrytao.wordpress.com/2026/09/24/were-gonna-need-a-lot-more-mathematicians/}.

\bibitem%[SAIR26]
{SAIR26} T. Tao, \emph{SAIR's open math model initiative}. What's new (blog), 18 September 2026. \url{https://terrytao.wordpress.com/2026/09/18/sairs-open-math-model-initiative/}.

\bibitem%[sch66]
{sch66} J. Schmidt, \emph{A general existence theorem on partial algebras and its special cases}. Colloq. Math. \textbf{14} (1966) pp.~73--87.

\bibitem%[sch70]
{sch70} J. Schmidt, \emph{A homomorphism theorem for partial algebras}. Colloq. Math. \textbf{21} (1970) pp.~5--21.

\bibitem%[Sch24]
{Sch24} M. Sch\"{o}nfinkel, \emph{\"{U}ber die Bausteine der mathematischen Logik}. Math. Ann. \textbf{92}(3) (1924) pp.~305--316.

%\bibitem%[Sch55]
%{Sch55} M.~P. Sch{\"{u}}tzenberger, \emph{Une th\'{e}orie alg\'{e}brique du codage}. S\'{e}minaire Dubreil. Alg\`{e}bre et th\'{e}orie des nombres, \textbf{9} (1955--1956), Expos\'{e} No. 15, pp.~1--24.

\bibitem%[Sch22]
{Sch22} P. Scholze, \emph{Liquid Tensor Experiment}. Experimental Mathematics \textbf{31} (2022), pp.~349--354.

\bibitem%[Sel06]
{Sel06} J.~P. Seldin, \emph{The logic of Curry and Church}. In: D.~M. Gabbay and J. Woods (Eds.) Handbook of the History of Logic. Vol.~\textbf{5}. Logic from Russell to Church (2009), pp.~819--873. Elsevier/North-Holland, Amsterdam, 2009.

\bibitem%[sl68]
{sl68} J. S{\l}omi\'{n}ski, \emph{Peano-algebras and quasi-algebras}.
Dissertationes Math. (Rozprawy Mat.) \textbf{57} (1968).

%\bibitem%[sl74]
%{sl74} J. S{\l}omi\'{n}ski, \emph{On the greatest congruence relation contained in an equivalence relation and its applications to the algebraic theory of machines}. Colloq. Math.  \textbf{29} (1974), pp.~31--43.

\bibitem{SOK94}
C.~C. Squier, F. Otto and Y. Kobayashi, \emph{A finiteness condition for rewriting systems}. Theoretical Computer Science \textbf{131} (1994), pp.~271--294.

\bibitem{Str76}
R. Street, \emph{Limits indexed by category-valued 2-functors}. Journal of Pure and Applied Algebra \textbf{8}.2 (1976), pp.~149-181.

\bibitem{Str87}
R. Street, \emph{The algebra of oriented simplexes}. Journal of Pure and Applied Algebra \textbf{49}.3 (1987), pp.~283--335.

%\bibitem%[Ste81]
%{Ste81} M. Steinby, \emph{Some algebraic aspects of recognizability and rationality}. In F. G\'{e}cseg, editor, \emph{Fundamentals of Computation Theory: Proceedings of the 1981 International FCT-Conference, Szeged, Hungary, August 24--28 (1981)} pp.~360--372, Berlin-Heidelberg.

\bibitem%[ST00]
{ST00} M. Steinby and W. Thomas, \emph{Trees and term rewriting in 1910: on a paper by Axel Thue}. Bull. Eur. Assoc. Theor. Comput. Sci. EATCS \textbf{72} (2000), pp.~256--269.

\bibitem%[St72]
{St72} S. Stenlund, \emph{Combinators, $\lambda$-terms and proof theory}. D. Reidel Publishing Co., Dordrecht, 1972.

\bibitem%[Sum26]
{Sum26} M. Abouzaid, M. Brenner, I. Corwin, S. Fomin, D. Freed, J. G\'{o}mez-Serrano, L. Guth, J. Harris, M. Kisin, B. Kra, D. Litt, M. Matchett Wood, C. McMullen, A. Naor, M. Olsson, T. Perutz, E. Riehl, R. Rouquier, T. Scanlon, G. Staffilani, R. Vakil, R. Ward, L. Williams and C. Xu, \emph{Report of Summit on PhD Math Education in the Age of AI}. Summit held on 17--18 September 2026. \url{https://cmsa.fas.harvard.edu/media/2026/09/Summit-on-PhD-Math-Education-in-the-Age-of-AI.pdf}.

\bibitem%[T08]
{T08} H. Tietze, \emph{\"{U}ber die topologischen Invarianten mehrdimensionaler Mannigfaltigkeiten}, Monat. Math. Phys. \textbf{19} (1908), pp.~1--118.

\bibitem%[TaoFM26]
{TaoFM26} T. Tao et al., \emph{A severe misalignment of AI in mathematics}. Joint declaration of 25 Fields medallists, What's new (blog), 11 September 2026. \url{https://terrytao.wordpress.com/2026/09/11/a-severe-misalignment-of-ai-in-mathematics/}.

\bibitem%[TBG91]
{tbg91} A. Tarlecki, R. Burstall and J. Goguen, \emph{Some fundamental algebraic tools for the semantics of computation: Part 3. Indexed categories}, Theoretical Computer Science, \textbf{91} (1991), pp.~239--264.

%\bibitem%[Tei51]
%{Tei51} M. Teissier, \emph{Sur les \'{e}quivalences r\'{e}guli\`{e}res dans les demi-groupes}. C.R. Acad. Sci. Paris \textbf{232} (1951), pp.~1987--1989.

\bibitem%[Tha69]
{Tha69} J.~W. Thatcher, \emph{Transformations and translations from the point of view of generalized finite automata theory}. In \emph{Proceedings of the ACM Symposium on Theory Computing} (1969),
pp.~129--142.

%\bibitem%[Tha73]
%{Tha73} J.~W. Thatcher, \emph{Tree automata: an informal survey}. In A.~V. Aho, editor, \emph{Currents in the Theory of Computing}, Prentice-Hall, Englewoods Cliffs, N.~J., pp.~143--172, 1973.

%\bibitem%[TW65]
%{TW65} J.~W. Thatcher and J.~B. Wright, \emph{Generalized finite automata}. Notices Amer. Math. Soc. \textbf{12} (1965), Abstract No 65T-649, p.~820.
%
%\bibitem%[TW68]
%{TW68} J.~W. Thatcher and J.~B. Wright, \emph{Generalized finite automata theory with an application to a decision problem of second order logic}. MST \textbf{2} (1968), pp.~57--81.

\bibitem%[Th10]
{Th10} A. Thue, \emph{Die L\"{o}sung eines Spezialfalles eines generallen logischen Problems}. Videnskabs-Selskabets Skrifter Kristiania. I. Math.-Nat. Klasse 
No.~\textbf{8} (1910).
%, pp.~1--17. 

\bibitem%[Th12]
{Th12} A. Thue, \emph{\"{U}ber die gegenseitige Lage gleicher Teile gewisser Zeichenreihen}. Videnskabs-Selskabets Skrifter Kristiania. I. Math.-Nat. Klasse 
No.~\textbf{9} (1912).
%, pp.~1--67.

\bibitem%[Th14]
{Th14} A. Thue, \emph{Probleme \"{u}ber Ver\"{a}nderungen von Zeichenreihen nach gegebenen Regeln}. Videnskabs-Selskabets Skrifter Kristiania. I. Math.-Nat. Klasse No.~\textbf{10} (1914).
%, pp.~1--27.

\bibitem%[Th77]
{Th77} A. Thue, \emph{Selected Mathematical Papers}. With an introduction by Carl Ludwig Siegel and a biography by Viggo Brun. Edited by Trygve Nagell, Atle Selberg, Sigmund Selberg, and Knut Thalberg. Universitetsforlaget, Oslo-Bergen-Troms\o, 1977.

\bibitem%[Tur36]
{Tur36} A.~M. Turing, \emph{On Computable Numbers, with an Application to the Entscheidungsproblem}. Proc. London Math. Soc. \textbf{42} (1936), pp.~230--265.

\bibitem%[Vag26]
{Vag26} M. Vago Marzal, \emph{A Formalization of Higher-Order Categories in Lean~4} (Version 0.1.0). Zenodo, 2026. \url{https://doi.org/10.5281/zenodo.20631672}.

\bibitem%[VagRepo]
{VagRepo} M. Vago Marzal, \emph{Higher category theory in Lean~4}. Code repository and documentation. \url{https://mariovagomarzal.github.io/higher_category_theory/}.

%\bibitem%[Voe11]
%{Voe11} V. Voevodsky, \emph{Univalent foundations of mathematics}. In International Workshop on Logic, Language, Information, and Computation, pp.~4--4. Springer, 2011. 

%\bibitem%[w21]
%{w21} F. Wagemann, \emph{Crossed modules}. De Gruyter, Berlin, 2021.

\bibitem%[XTX26]
{XTX26} XTX Markets, \emph{2026 update on XTX Markets' AI philanthropy}. 16 September 2026. \url{https://www.xtxmarkets.com/news/2026-update-on-xtx-markets-ai-philanthropy/}.

\bibitem%[w92]
{w92} W. Wechler, \emph{Universal algebra for computer scientists}. Springer-Verlag, Berlin, 1992.

\bibitem%[Wol26]
{Wol26} S. Wolfram, \emph{What's the future for pure math research in the age of AI?} Stephen Wolfram Writings, 28 September 2026. \url{https://writings.stephenwolfram.com/2026/09/whats-the-future-for-pure-math-research-in-the-age-of-ai/}.

\bibitem%[uw90]
{uw90} U. Wolter, \emph{An algebraic approach to deduction in equational partial Horn theories}. Journal of Information Processing and Cybernetics. EIK, \textbf{27}, pp.~85--128, 1990.lag, Berlin, 1992.

\bibitem%[Yan26]
{Yan26} K. Yang, \emph{What happens when formalization becomes cheap?} Blog post, 22 September 2026. \url{https://yangky11.github.io/blog/ai-and-formalization/}.

%\bibitem%[y60]
%{y60} N. Yoneda, \emph{On Ext and exact sequences}. J. Fac. Sci. Univ. Tokyo Sect. I \textbf{8} (1960), pp.~ 507--576.


\end{thebibliography}

